\section{אפקט אהרונוב-בוהם (\LR{The Aharonov-Bohm Effect})}

אפקט זה מדגים את ההבדל המהותי בין מכניקה קלאסית לקוונטית ביחס למעמד הפוטנציאלים האלקטרומגנטיים.

\subsection{תיאור המערכת}
נתבונן בסליל (\LR{Solenoid}) אינסופי ברדיוס $a$ הנושא זרם $I$.
\begin{itemize}
    \item **בתוך הסליל ($r < a$):** קיים שדה מגנטי אחיד $B_z = B_0$.
    \item **מחוץ לסליל ($r > a$):** השדה המגנטי מתאפס $\vec{B} = 0$.
    \item **החלקיק:** חלקיק קוונטי (מטען $q$, מסה $m$) נע על טבעת ברדיוס $b > a$ המקיפה את הסליל.
\end{itemize}

באופן קלאסי, כוח לורנץ תלוי רק בשדות ($\vec{F} = q\vec{v} \times \vec{B}$). מכיוון שהחלקיק נע באזור בו $\vec{B}=0$, הוא לא אמור להרגיש שום כוח ולא להיות מושפע מהזרם בסליל. קוונטית, המצב שונה.

\subsection{חישוב הפוטנציאל הווקטורי}
במכניקת הקוונטים, ההמילטוניאן תלוי בפוטנציאלים ($\vec{A}, \phi$) ולא ישירות בשדות. נחשב את $\vec{A}$ בקואורדינטות גליליות תוך שימוש בסימטריה ($\vec{A} = A_\theta(r)\hat{\theta}$) ובחוק סטוקס:
\[
\oint \vec{A} \cdot d\vec{l} = \int (\nabla \times \vec{A}) \cdot d\vec{S} = \Phi_B
\]
כאשר $\Phi_B$ הוא השטף המגנטי. עבור מסלול ברדיוס $r$:
\[
A_\theta \cdot 2\pi r = \Phi_{enclosed}
\]

\begin{resultbox}{הפוטנציאל הווקטורי}
\begin{equation}
A_\theta(r) = \begin{cases} 
\frac{B_0 r}{2} & r < a \\
\frac{B_0 a^2}{2r} = \frac{\Phi_B}{2\pi r} & r > a 
\end{cases}
\end{equation}
מחוץ לסליל, למרות ש-$\vec{B}=0$, הפוטנציאל $\vec{A} \neq 0$.
\end{resultbox}

\subsection{פתרון משוואת שרדינגר על הטבעת}
ההמילטוניאן לחלקיק טעון הוא:
\[
H = \frac{1}{2m}(\vec{p} - q\vec{A})^2 + q\phi
\]
נניח $\phi=0$. התנועה מוגבלת למעגל ($r=b$), ולכן אופרטור התנע הוא משיקי: $p_\theta = -i\hbar \frac{1}{b}\frac{\partial}{\partial \theta}$.
נציב את הפוטנציאל שמצאנו:
\[
H = \frac{1}{2m} \left( -i\hbar \frac{1}{b}\frac{\partial}{\partial \theta} - q \frac{\Phi_B}{2\pi b} \right)^2
\]
נחפש פונקציות עצמיות מהצורה $\psi(\theta) = N e^{in\theta}$. תנאי שפה מחזוריים $\psi(\theta + 2\pi) = \psi(\theta)$ מחייבים $n \in \mathbb{Z}$.

\begin{resultbox}{ספקטרום האנרגיה}
\begin{equation}
E_n = \frac{\hbar^2}{2m b^2} \left( n - \frac{q\Phi_B}{h} \right)^2 = \frac{\hbar^2}{2m b^2} \left( n - \frac{\Phi_B}{\Phi_0} \right)^2
\end{equation}
כאשר $\Phi_0 = h/q$ הוא קוונט השטף.
\end{resultbox}

\textbf{משמעות פיזיקלית:}
\begin{enumerate}
    \item ללא שדה מגנטי ($\Phi_B=0$), הרמות $n$ ו-$-n$ מנוונות ($E \propto n^2$).
    \item נוכחות השטף מסירה את הניוון: האנרגיה משתנה למרות שהחלקיק לא עובר באזור בו יש שדה מגנטי.
\end{enumerate}

\subsection{פאזה גיאומטרית וניסוי התאבכות}
הניסוי המפורסם יותר של אהרונוב-בוהם הוא ניסוי התאבכות (בדומה לשני סדקים) כאשר הסליל ממוקם בין שני המסלולים האפשריים של האלקטרון.

פונקציית הגל בנוכחות פוטנציאל וקטורי צוברת פאזה ביחס למקרה החופשי ($\psi_0$):
\[
\psi = \psi_0 \exp\left( \frac{iq}{\hbar} \int \vec{A} \cdot d\vec{l} \right)
\]
הפרש הפאזה בין שני המסלולים (1 ו-2) הוא:
\begin{equation}
\Delta \varphi = \frac{q}{\hbar} \left( \int_1 \vec{A}\cdot d\vec{l} - \int_2 \vec{A}\cdot d\vec{l} \right) = \frac{q}{\hbar} \oint \vec{A} \cdot d\vec{l} = \frac{q\Phi_B}{\hbar}
\end{equation}
הפרש פאזה זה גורם להסטת תבנית ההתאבכות על המסך.

\begin{theorembox}{מסקנה קונספטואלית: מעמד הפוטנציאלים}
במכניקה קלאסית, $\vec{A}$ הוא כלי עזר מתמטי התלוי בכיול (\LR{Gauge}). במכניקת הקוונטים, הפוטנציאלים הם הגדלים היסודיים המופיעים במשוואות.
עם זאת, התוצאה הפיזיקלית (האנרגיה או הסטת הפאזה) תלויה ב-$\Phi_B = \oint \vec{A} \cdot d\vec{l}$, שהוא גודל אינווריאנטי לכיול.
\end{theorembox}

\section{מבוא לספין (\LR{Spin})}

המוטיבציה למושג הספין הגיעה מניסויים ספקטרוסקופיים (כגון המבנה העדין ואפקט זימן האנומלי) שהראו פיצול ברמות האנרגיה שלא תאם את התחזית של תנע זוויתי אורביטלי $L$.

\subsection{השערת אולנבק וגאודסמיט והפרדוקס הקלאסי}
הוצע שלאלקטרון יש תנע זוויתי פנימי $S$ עם מספר קוונטי $s=1/2$, כך שהניוון הוא $2s+1=2$.
האם ניתן לתאר זאת ככדור מסתובב?
\begin{examplebox}{הפרדוקס הקלאסי}
נניח שהאלקטרון הוא כדור ברדיוס $r$ המסתובב עם תנע זוויתי $L \approx \hbar/2$.
\[
L \sim mvr \approx \hbar \implies v \approx \frac{\hbar}{mr}
\]
אם נציב את הרדיוס הקלאסי של האלקטרון (או חסם עליון ניסיוני של $10^{-15}$ מטר), נקבל שמהירות הסיבוב הנדרשת היא:
\[
v \gg c
\]
מסקנה: הספין אינו סיבוב פיזי של המסה במרחב, אלא דרגת חופש פנימית יסודית.
\end{examplebox}

\subsection{סימטריה וחבורת פואנקרה}
הספין נובע מתכונות הסימטריה של המרחב-זמן (חבורת פואנקרה). חלקיקים אלמנטריים מסווגים לפי המסה והספין שלהם.
תכונה ייחודית של ספין 1/2 היא שסיבוב של $2\pi$ במרחב אינו מחזיר את פונקציית הגל לעצמה, אלא הופך סימן ($\psi \to -\psi$). נדרש סיבוב של $4\pi$ כדי לחזור למצב המקורי.

\subsection{הפורמליזם המתמטי: חבורת $SU(2)$}
מערכת עם שתי רמות מתוארת על ידי וקטור במרחב דו-ממדי (ספינור): $\chi = \begin{pmatrix} \alpha \\ \beta \end{pmatrix}$.
הטרנספורמציות ששומרות על הנורמה במרחב זה שייכות לחבורה $SU(2)$ (מטריצות אוניטריות $2\times 2$ עם דטרמיננטה 1).

כל איבר בחבורה ניתן לכתיבה בצורה:
\[
U = e^{iX}
\]
כאשר $X$ היא מטריצה הרמיטית. הדרישה $\det(U)=1$ גוררת תנאי על העקבה (\LR{Trace}) של $X$:
\[
\det(e^{iX}) = e^{i \text{Tr}(X)} = 1 \implies \text{Tr}(X) = 0
\]
.

\subsubsection{מטריצות פאולי}
מטריצה הרמיטית כללית $2\times 2$ עם עקבה 0 תלויה ב-3 פרמטרים ממשיים וניתנת לפריסה על ידי 3 מטריצות בסיס - **מטריצות פאולי**:
\begin{definitionbox}{מטריצות פאולי ($\vec{\sigma}$)}
\begin{equation}
\sigma_x = \begin{pmatrix} 0 & 1 \\ 1 & 0 \end{pmatrix}, \quad
\sigma_y = \begin{pmatrix} 0 & -i \\ i & 0 \end{pmatrix}, \quad
\sigma_z = \begin{pmatrix} 1 & 0 \\ 0 & -1 \end{pmatrix}
\end{equation}
אופרטור הספין מוגדר כ: $\vec{S} = \frac{\hbar}{2}\vec{\sigma}$.
\end{definitionbox}